\documentclass[10pt,a4paper]{amsart}
\usepackage[margin=19mm]{geometry}
\usepackage{amsmath,amssymb,mathtools}
\usepackage[scr=rsfs,cal=euler]{mathalfa}
\usepackage{xfrac}
\usepackage[unicode,hidelinks,bookmarks=false]{hyperref}
\newcommand{\ie}{i.e.\ }
\newcommand{\cf}{cf.\ }
\newcommand{\resp}{respectively\ }
\newcommand{\Subsection}{\subsection}
\providecommand{\nobreakdash}{\nobreak\hbox{-}}
\setlength{\emergencystretch}{2em}
\multlinegap0pt
\usepackage{bm}
\usepackage{bbm}
\usepackage{dsfont}
\usepackage{xspace}
\usepackage{cancel}
\usepackage{mathtools}
\usepackage{amsthm}
\makeatletter
\@ifundefined{Theorem}{\newtheorem{Theorem}{Theorem}}{}
\@ifundefined{Lemma}{\newtheorem{Lemma}[Theorem]{Lemma}}{}
\makeatother
\title[Regularized $\zeta\_{\Delta}(1)$ for Polyhedra]{Regularized $\zeta\_{\Delta}(1)$ for Polyhedra}
\author{Alexey Yu. Kokotov, Dmitrii V. Korikov}
\date{Source: arXiv:2502.03351v3 (CC BY 4.0)}
\begin{document}
\maketitle

\noindent\emph{Source and licence.} Regularized $\zeta_{\Delta}(1)$ for Polyhedra. Version arXiv:2502.03351v3 (CC BY 4.0), distributed under the Creative Commons Attribution 4.0 International licence, as recorded in the arXiv licence metadata for this exact version and shown on the abstract page https://arxiv.org/abs/2502.03351v3. Full rights notice: licenses/arxiv-2502.03351-CC-BY-4.0.txt.\par\smallskip

\noindent\textit{Editorial scope.} Counted unit: the Lemma whose source label is \texttt{Yamada lemma} in the article (source0.tex lines 803--828, source label \texttt{Yamada lemma}), delivered together with the article's own complete original proof of it (source0.tex lines 829--859, one \texttt{proof} environment of 31 lines). No mathematics is written, repaired or re-derived: statement and proof are the article's own text line for line. Required context retained uncounted: none. Every definition, display and label of those lines is retained unchanged. Omitted from this package and not redistributed: 1--802, 860--1425. Nothing inside a retained excerpt is omitted or altered: every retained body is delivered exactly as the article's lines are written, so no omission receipt of any kind accompanies this package. Citations: the retained excerpts carry no citation command at all, so no bibliography entry of the article is transcribed and no reference is redistributed. Reference adaptation: none was needed and none was made; every reference inside the delivered unit resolves inside it, so no unresolved reference or citation is produced. Printed slips kept literal: no printed word, symbol, hypothesis or number of the counted statement or of its proof was altered. Layout and class: the article's own class is \texttt{sigma} with packages including ifpdf, mathrsfs; the wrapper is the lane's pinned \texttt{amsart} editorial wrapper at 10pt on A4, which restates the theorem-like environments the retained text needs and adds the article's own macro definitions that the retained text needs (none), with extra wrapper packages (bm, bbm, dsfont, xspace, cancel, mathtools). The delivered numbering is the unit's own. Assets: no retained span contains a figure environment, a graphics-inclusion command, a PDF-page inclusion, a table or a TikZ picture; no image file is copied into this package, the package declares no asset dependency and no image file is redistributed with it. Licence: version arXiv:2502.03351v3 of this article is distributed under the Creative Commons Attribution 4.0 International licence, as recorded in the arXiv licence metadata for this exact version and shown on the abstract page https://arxiv.org/abs/2502.03351v3. A full rights notice is delivered at licenses/arxiv-2502.03351-CC-BY-4.0.txt. Characters: every retained excerpt is pure ASCII, so the delivered engine, XeLaTeX with its default Latin Modern text font, reports no missing character.
\begin{Lemma}\label{Yamada lemma}\qquad
\begin{enumerate}\itemsep=0pt
\item[$(i)$] Let $v_\pm$ be holomorphic differentials on $X_\pm$, respectively. Then there is a family $v=v(\varepsilon)$ of holomorphic differentials on $X=X(\varepsilon)$ obeying the inequality
\begin{equation}
\label{Yamada est}
\sum_{\pm}\|v-v_\pm\|_{L_2(X_\pm\backslash U)}+\|v\|_{L_2(U)}\le c\varepsilon,
\end{equation}
where $U=U(\varepsilon)$ is given by
$U:=\{P\in X | |\xi(P)\pm 1|\ge c\}$
and the $L_2$-norm of a one-form~$v$ on subdomain $U\subset X$ is given by
\[\|v\|_{L_2(U)}:=\left(\int_{U}v\wedge\overline{\star v}\right)^{1/2}.\]
\item[$(ii)$] Let $W$ and $W_\pm$ be the canonical bimeromorphic differentials on $X$ and $X_\pm$, respectively. Then the inequality
\begin{equation}
\label{Yamada est merom}
\|W(\cdot,y)-W_\pm(\cdot,y)\|_{L_2(X_\pm\backslash U)}+\|W(\cdot,y)\|_{L_2(X_\mp\cup U)}\le \frac{c\varepsilon}{ d(y,P_\pm)^2}
\end{equation}
holds for any $y\in X_\pm\backslash U$.
\item[$(iii)$] Denote $W_0=(\xi-\xi')^{-2}{\rm d}\xi {\rm d}\xi'$. Then there is a meromorphic differential $Y=Y(\cdot,\xi')$ on~$X$ obeying
\begin{equation}
\label{Yamada est on pinch}
\|Y(\cdot,\xi')\|_{L_2(X\backslash V)}+\|W_0(\cdot,\xi')-Y(\cdot,\xi')\|_{L_2(V)}\le \frac{c}{\big|1-\xi'^2\big|},
\end{equation}
where $\xi'\in V$ and $V$ us given by
$V:=\{P\in X | |z_+(P)|\le c \text{ or } |z_-(P)|\le c\}$.
\end{enumerate}
\end{Lemma}

\begin{proof}
$(i)$ We have
\begin{equation}
\label{Yamada y}
y_{\pm}(P):=\int_{P_\pm}^P v_\pm=\sum_{k>0}c_{k}z_\pm^{k}=\sum_{k>0}\frac{c_{k}\varepsilon^{k}}{\bigl(1-\xi^2\bigr)^k},
\end{equation}
where $z_\pm=z_\pm(P)$, $\xi=\xi(P)$, and $P\in \partial U$.
Let $h=h(\varepsilon)$ be the harmonic extension of $y_\pm|_{\partial U}$ into $U$; then the above expansion implies $\|{\rm d}h\|_{L_{2}(U)}\le c\varepsilon$. Introduce the form $\Phi=\Phi(\varepsilon)$ given by $\Phi=v_\pm$ on $X_\pm\backslash U$ and by $\Phi={\rm d}h$ on $U$. By construction, the jump of the form $\Phi$ on $\partial U$ is zero on any vector field $k$ tangent to $\partial U$. Due to this, ${\rm d}\Phi$ has no singularities on $\partial U$; moreover, we have ${\rm d}\Phi=0$ on the whole $X$. Note that the form $\tilde{\Phi}:=\Phi-{\rm i}\star\Phi$ vanishes outside $U$ due to~${{\rm i}\star {\rm d}z={\rm d}z}$; as a corollary, we have
\begin{equation}
\label{Yamada alpha est h}
\|\tilde{\Phi}\|_{L_2(X)}\le 2\|{\rm d}h\|_{L_2(X)}\le c\varepsilon.
\end{equation}
Consider the Hodge orthogonal decomposition $\tilde{\Phi}={\rm d}\alpha+\delta\beta+\gamma$, where ${\rm d}\gamma=0$, $\delta\gamma=0$ and
\begin{equation}
\label{Yamada alpha est}
\|{\rm d}\alpha\|_{L_2(X)}\le\|\tilde{\Phi}\|_{L_2(X)}\le c\varepsilon.
\end{equation}
Then we have
$\Phi-{\rm d}\alpha=\delta\beta+\gamma+{\rm i}\star\Phi$,
where the left-hand side is closed and the right-hand side is coclosed. So, the form $\Phi-{\rm d}\alpha$ is harmonic and, thus, the form
$v=\frac{1}{2}\bigl(\Phi-{\rm d}\alpha+{\rm i}\star(\Phi-{\rm d}\alpha)\bigr)$
is holomorphic. Note that
\[v-v_\pm=\frac{1}{2}({\rm d}\alpha+{\rm i}\star {\rm d}\alpha) \qquad\text{on}\ X_\pm\backslash U, \qquad v=\frac{1}{2}( d(h-\alpha)+{\rm i}\star d(h-\alpha)) \qquad\text{on}\ U.\]
Then \eqref{Yamada alpha est h}, \eqref{Yamada alpha est} imply \eqref{Yamada est}.

$(ii)$ Repeating the above construction for $v_\pm=W_\pm(\cdot,y)$ and $v_\mp=0$, one gets a differential $v$ obeying \eqref{Yamada est}, where the additional multiplier $ d(y,P_\pm)^{-2}$ should appear in the right-hand side due to the standard estimate $|c_k|\le c|z_\pm(y)|^{-(k+2)}$ for the coefficients in \eqref{Yamada y}. By construction, $v$ and $W(\cdot,y)$ share the same singularity on $X$, i.e., $v-W(\cdot,y)$ is holomorphic on $X$. In addition, the periods of $v$ coincide with those of $(1/2){\rm i}\star {\rm d}\alpha$. Since $\tilde{\Phi}={\rm d}\alpha+\delta\beta+\gamma$ vanishes outside $U$, ${\rm d}\alpha$ is harmonic. Then the increasing smoothness theorems for solutions to elliptic equations imply
\[\Big|\int_l\star {\rm d}\alpha\Big|\le |l|\|{\rm d}\alpha\|_{L_2(X)}\le c|l|\varepsilon d(y,P_\pm)^{-2},\]
where $l$ is any path which does not intersect $V$ and $|l|$ is its length. Due to the last estimate (with $l=a_{k,\pm},b_{k,\pm}$) and the Riemann bilinear relations, we have $v-W(\cdot,y)=\tilde{v}$, where $\|\tilde{v}\|_{L_2(X)}\le c|l|\varepsilon d(y,P_\pm)^{-2}$. As a corollary, we obtain~\eqref{Yamada est merom}.

$(iii)$ To prove $(iii)$, it is sufficient to repeat the above construction with the following replacements: (a)~instead of \eqref{Yamada y}, we consider the anti-derivative $y(P)=-(\xi-\xi')^{-1}$ of $W_0(\cdot,\xi')$ in the domain $V$, (b)~$h$~is defined as the harmonic extension of $y|_{\partial V}$ into domains $X\backslash V=(X_+\backslash V)\cup (X_-\backslash V)$, (c)~$\Phi$~is defined by $\Phi=W_0$ on $V$ and by $\Phi={\rm d}h$ on $X\backslash V$.
\end{proof}

\end{document}
